\documentclass[11pt]{article}
\usepackage[margin=1in]{geometry}
\usepackage{amsmath,amssymb,booktabs,hyperref}
\newcommand{\dd}{\mathrm d}
\title{Finite-time collective response and soft particle matching}
\author{Zimmerman programme: CPB-PR-2026-10-16}
\date{16 October 2026}
\begin{document}
\maketitle
\section{The question and existing obstructions}
The preceding checkpoint found a tiny direct Higgs-induced response at
chosen galactic densities. Could a stable medium nevertheless amplify it
through near-critical susceptibility? This report supplies a finite-time
test of that proposal, and calculates a previously omitted soft-Higgs loop
contribution to the daughter coupling. Neither supplies a completed galaxy
mechanism or a microscopic origin for the vacuum energy.

Earlier project work already derives the protected carrier's dilute coherent
phase pressure proportional to the square of its chemical-potential variable,
and shows why simply assigning it a MOND phonon is unjustified. A separate
barotropic-pressure argument already excludes using a single density-only
equation of state to support different exactly flat speeds over overlapping
density intervals. These results are retained, not rediscovered here.
The superfluid dark-matter route is also established literature:
Ref.~\cite{pr:superfluid} explicitly assumes its phonon action and a baryon
coupling; it does not derive those ingredients from this scalar carrier.

\section{A finite-time bound for a stable quadratic sector}
Consider physical coordinates $q$ with constant real positive-definite
kinetic matrix $M$ and positive-semidefinite stiffness $K$:
\begin{equation}
 L={1\over2}\dot q^TM\dot q-{1\over2}q^TKq
                +J(t)b^Tq,\qquad q(0)=\dot q(0)=0.
 \label{pr:system}
\end{equation}
Let $U$ orthogonally diagonalize $M^{-1/2}KM^{-1/2}$, with eigenvalues
$\Omega_j^2\geq0$. For a unit constant source the reciprocal observable
$b^Tq$ has step response
\begin{equation}
 R_b(t)=\sum_j f_j{1-\cos(\Omega_jt)\over\Omega_j^2},\qquad
 f_j=\bigl(U_j^TM^{-1/2}b\bigr)^2\geq0,\quad
 \sum_j f_j=b^TM^{-1}b.
 \label{pr:step}
\end{equation}
The zero-frequency summand is defined by its limit $f_jt^2/2$.
Since $0\leq1-\cos x\leq x^2/2$,
\begin{equation}
 0\leq R_b(t)\leq {t^2\over2}b^TM^{-1}b.
 \label{pr:bound}
\end{equation}
This proves the bound for arbitrary finite dimension, rather than inferring
it from numerical sampling. The static susceptibility, when it exists, is
$\chi_b(0)=b^TK^{-1}b=\sum_j f_j/\Omega_j^2$. It can diverge as a
mode softens while the finite-time bound remains unchanged.
Reaching a fraction $f>0$ of that static response requires
\begin{equation}
 t\geq\sqrt{{2f\chi_b(0)\over b^TM^{-1}b}}.
 \label{pr:minimumtime}
\end{equation}
This is necessary, not sufficient, and does not establish settling.
Undamped modes may overshoot their static response and need never settle.

The result also has a useful cross-observable version. For an observable
$a^Tq$ and any source satisfying $|J(s)|\leq J_{\max}$ on $0\leq s\leq t$,
the retarded mode kernel is $\sin(\Omega_j(t-s))/\Omega_j$.
Its absolute value is at most $t-s$. Cauchy--Schwarz on the mode
coefficients therefore gives
\begin{equation}
 |a^Tq(t)|\leq {J_{\max}t^2\over2}
       \sqrt{(a^TM^{-1}a)(b^TM^{-1}b)}.
 \label{pr:cross}
\end{equation}
This proof does not require reciprocal positive spectral weights for the
cross response. A regulated positive oscillator bath can be included among
the coordinates before elimination. A continuum extension requires finite
inertial sums and control of its limit; an unregulated effective kernel
does not automatically meet these hypotheses.

The theorem assumes zero perturbation initial data, a fixed stable
quadratic expansion, canonical second-order dynamics and a bounded source.
Preloaded excitations, a changing background, nonlinear phase transitions,
active gain or unstable modes are not covered. Their energy source and
evolution must instead be calculated. Neither a universal halo exclusion
nor a new general theorem of mathematical physics is claimed.

\section{Application to the existing scalar--polarization construction}
For a single real Fourier component, use the previously proposed static
energy and bare kinetic matrices
\begin{equation}
 K_k=\begin{pmatrix}Ck^2&-k\\-k&C^{-1}+\Pi_k\end{pmatrix},\qquad
 M=\begin{pmatrix}Z_\varphi&0\\0&I_P\end{pmatrix},\qquad
 C,Z_\varphi,I_P>0,\quad k>0,\quad\Pi_k>0.
\end{equation}
The determinant is $Ck^2\Pi_k>0$. A source on $\varphi$ has
\begin{equation}
 \chi_{\varphi\varphi}(0,k)
 ={C^{-1}+\Pi_k\over Ck^2\Pi_k}
 ={1\over Ck^2}+{1\over C^2k^2\Pi_k},\qquad
 R_{\varphi\varphi}(t,k)\leq {t^2\over2Z_\varphi}.
 \label{pr:twofield}
\end{equation}
Reducing $\Pi_k$ increases the static response but does not remove the
necessary time in Eq.~\eqref{pr:minimumtime}. For the dimensionless example
$C=Z_\varphi=I_P=k=1$, $\Pi_k=10^{-6}$, the static susceptibility is
$1000001$, while the exact step response at $t=10$ is only $25.25115$.
Reaching half the static value cannot occur before $t=1000.0005$ by the
bound. These are dimensionless controls, not galaxy times in years.

This two-coordinate numerical example freezes $\Pi_k$ as an instantaneous
stiffness. The actual nodal-fermion construction has a frequency-dependent
retarded kernel, positive absorption and induced inertia. Its earlier
positive-bath representation allows the theorem to be applied to a
regulated full quadratic realization when the stated positivity and finite
inertial-sum assumptions hold. The frozen two-coordinate curve is not
promoted to that fermionic medium's actual formation history.

The prior isotropic gapless branch has $\Pi_k\simeq\eta k$ statically
and $\Omega_{\rm slow}^2\simeq C^2\eta k^3/Z_\varphi$ at small $k$.
The static susceptibility scales as $k^{-3}$ in that regime; its response
time grows correspondingly as $k^{-3/2}$ at least. Actual damping and
settling are further requirements. None of the unfixed dimensional kinetic
coefficients has been calibrated by the dimensionless calculation here.

There is a dimensional necessary condition if this same bare scalar sector
is required to have propagation speed $C/Z_\varphi\leq c^2$. Define its
static enhancement over its own uncoupled scalar response by
$\mathcal A=Ck^2\chi_{\varphi\varphi}(0,k)
=1+(C\Pi_k)^{-1}$. Equations~\eqref{pr:minimumtime} and
\eqref{pr:twofield}, with $L=1/k$, imply
\begin{equation}
 t\geq {L\over c}\sqrt{2f\mathcal A},\qquad
 \mathcal A\leq {1\over2f}\left({ct\over L}\right)^2
 \quad\hbox{if fraction $f$ is reached by time $t$}.
 \label{pr:causalbound}
\end{equation}
For half-response in 10 Gyr at $L=1$ kpc, the necessary enhancement ceiling
is about $9.4\,10^{12}$. A putative $10^{47}$ enhancement in this same
channel would take at least $1.03\,10^{27}$ years to reach half-response
under these assumptions. This does not identify the earlier contact shift
with $\mathcal A$; matching their source and observable is still required.
It closes the specific shortcut of freely increasing static susceptibility
while retaining arbitrarily fast, initially unperturbed subluminal response.
It does not exclude a nonlinear or preloaded state, or a different source
normalization. The additional bare-speed assumption is not proved for the
current preferred-frame band model or for an unspecified ultraviolet theory.

The deterministic run checks 48 two-field parameter combinations and six
three-coordinate cross responses. Generalized-eigenmode responses agree
with independently integrated ODEs in the latter cases. The resulting 648
controls test implementations and examples; the displayed inequalities are
proved by the spectral argument above.

\section{A normal-gas comparison with an explicit time normalization}
As a separate controlled surrogate, take a homogeneous mean-subtracted
Newtonian fluid, fixed density $\rho_0$, adiabatic sound speed $c_s$, and
no expansion. Linear continuity, Euler and Poisson equations give
\begin{equation}
 \ddot{\delta\rho_c}+(c_s^2k^2-\omega_J^2)\delta\rho_c
       =\omega_J^2\delta\rho_b,
 \qquad \omega_J^2=4\pi G\rho_0.
\end{equation}
For $k\geq k_J=\omega_J/c_s$, zero initial perturbation and a constant
baryon-density Fourier source, define $a=k/k_J$ and $\tau=\omega_Jt$.
The density ratio is
\begin{equation}
 {\delta\rho_c(t)\over\delta\rho_b}
 ={1-\cos(\sqrt{a^2-1}\,\tau)\over a^2-1}
 \leq {\tau^2\over2}.
 \label{pr:gas}
\end{equation}
At $a=1$ the limiting equality holds; there is no finite static solution.
At $a>1$ the static ratio is $(a^2-1)^{-1}$. Below $k_J$ the system is
Jeans unstable and this stable-mode argument is not applicable.

For the selected comparison $\rho_0=0.3\,\mathrm{GeV\,cm}^{-3}/c^2$
and one-dimensional dispersion $150\,\mathrm{km\,s}^{-1}$, let
$c_s^2=(5/3)\sigma_v^2$. Then
\begin{equation}
 \omega_J^{-1}=47.3144\,\mathrm{Myr},\qquad
 k_J^{-1}=9.37049\,\mathrm{kpc},\qquad
 {1\over2}(\omega_J\,10\,\mathrm{Gyr})^2=22334.9.
\end{equation}
Here $k_J^{-1}$ is an inverse wavenumber, not the wavelength $2\pi/k_J$.
For example $a=1.000001$ has static ratio about $499999.75$ but a
10-Gyr step ratio about $22169.10$. Critical static response does not
yield arbitrarily rapid amplification.

This is a linear fluid comparison, not the earlier collisionless Maxwellian
Volterra calculation, nor a physical finite halo equilibrium. Large density
contrasts would invalidate linearization; the inequalities can be tested
with arbitrarily small source amplitudes. These density-response units
must not be equated to the earlier dimensionless contact shift: connecting
the contact to the actual source vector $b$ and force observable $a$ is a
missing matching calculation. The value 22334.9 is therefore not a
universal bound on the theory's galactic amplification.

\section{The soft-Higgs daughter coupling is not exactly zero}
The selected scalar action has $g(X^*\Phi^2+X\Phi^{*2})$ and
$\lambda_{HX}|H|^2|X|^2$ in its potential. Even with a zero daughter
portal at tree matching, a cubic loop with an internal daughter and source
produces a finite soft-Higgs vertex. Let the external Higgs momentum vanish
and the external daughter be on shell, $p^2=m^2$.

The normalization and sign can be checked directly in the Euclidean
functional integral. Integrating out the quadratic complex $X$ field gives
the nonlocal action term $-g^2\Phi^{*2}G_X\Phi^2$. Contracting one
daughter pair gives a quadratic correction $-4g^2G_XG_\Phi$; the factor
four comes from the two choices at each end. Its mass-squared correction,
up to an $M$-independent subtraction at fixed other parameters, is
\begin{equation}
 \Sigma_{\rm mass}(p^2)={g^2\over4\pi^2}
 \int_0^1\dd x\,\ln{D(x,p^2)\over\mu_R^2}+\text{constant},\quad
 D=xM^2+(1-x)m^2-x(1-x)p^2.
\end{equation}
The divergent subtraction in the cubic mass diagram does not affect its
finite derivative with respect to $M^2$. Inserting the soft Higgs through
$\delta M^2=\lambda_{HX}vh$ therefore gives the leading coefficient
of $-vh|\Phi|^2$ in the Lagrangian:
\begin{equation}
 C_{H\Phi}^{\rm soft,loop}
 =\lambda_{HX}{g^2\over4\pi^2}
 \int_0^1{ x\,\dd x\over xM^2+(1-x)^2m^2}>0.
 \label{pr:loop}
\end{equation}
External-leg corrections enter beyond this leading order when the tree
daughter portal is zero. A symmetry-allowed finite daughter-portal matching
term can add to this result; the loop contribution alone is not a
renormalization-condition-independent prediction for the full coupling.
The statement is that the specified tree zero is not an exact vanishing
of this perturbative soft vertex.
This vertex respects the dark global $U(1)$. That symmetry statement alone
does not extend a classical coldness bound to the full renormalized action.

As an analytic normalization control, at $M=2m$ the integral is
\begin{equation}
 {1\over m^2}\int_0^1{x\over(1+x)^2}\dd x
 ={\ln2-1/2\over m^2}.
\end{equation}
Independent quadratures and a mass-derivative check are performed at six
mass ratios. The selected 2-MeV candidate gives approximately
\begin{equation}
 {C_{H\Phi}^{\rm soft,loop}\over\lambda_{HX}}
 \simeq6.50\,10^{-8},\qquad
 C_{H\Phi}^{\rm soft,loop}\simeq1.02\,10^{-17}.
\end{equation}
With the same slow nonrelativistic nucleon medium as the previous
checkpoint, $vh=-f_N\rho_b/m_h^2$. Thus this loop contribution gives
\begin{equation}
 \delta m_\Phi^2=-C_{H\Phi}^{\rm soft,loop}{f_N\rho_b\over m_h^2}.
\end{equation}
At the chosen $0.1\,M_\odot\,\mathrm{pc}^{-3}$ density its magnitude is
of order $6\,10^{-63}\,\mathrm{GeV}^2$, a fractional daughter
mass-squared shift of order $1.5\,10^{-57}$. This perturbative soft
effect does not repair the galaxy-response mechanism. Its smallness is
not a complete proof about collective effects or a different finite
matching coupling. Equation~\eqref{pr:loop} also cannot be used as the
on-shell $125$-GeV Higgs decay form factor: the soft and hard kinematics
are different. No new freeze-in abundance is inferred from it.

\section{What changed and what remains}
The static critical construction now has a necessary finite-time response
condition independent of how small its stiffness is made. The particle
bridge now includes an explicit leading radiative soft coupling, with its
normalization, sign and matching caveat. These are constraints on the
route to closure, not closure itself.

A successful continuation must calculate the physical source vector and
kinetic normalization of one occupied response state, apply the time bound
with those units, and solve the nonlinear formation and energy receiver.
The nodal medium remains a separately postulated occupied sector; the
protected scalar carrier is not automatically that sector. Their bridge,
source selection, finite mass supply, actual metric/lensing and cosmological
tests remain required. The additive vacuum magnitude and the normalization
of $a_0=\kappa c\sqrt{G\rho_{\rm DE}}$ remain unfixed. Internal scalar
number-changing reactions still have zero vacuum energy-transfer $Q$.

Evidence is in \texttt{passive\_response\_run\_001} (648 bounded controls)
and \texttt{soft\_higgs\_matching\_run\_001} (19 controls), with
\texttt{causal\_response\_run\_001} (three unit controls).
The mathematical and physical review is self-review. Neither a globally
novel mechanism nor independent referee
certification is claimed.

\begin{thebibliography}{9}
\bibitem{pr:superfluid} L. Berezhiani and J. Khoury,
Theory of Dark Matter Superfluidity,
\href{https://arxiv.org/html/1507.01019v2}{arXiv:1507.01019v2},
Sec.~1.3 and Sec.~4. Cited only as established context for a distinct
phonon construction; no current observational verdict is imported.
\end{thebibliography}
\end{document}
